\PassOptionsToPackage{unicode}{hyperref}
\PassOptionsToPackage{hyphens}{url}
\PassOptionsToPackage{dvipsnames,svgnames,x11names}{xcolor}
\documentclass[
  11pt,
]{article}
\usepackage{amsmath,amssymb}
\usepackage{iftex}
\ifPDFTeX
  \usepackage[T1]{fontenc}
  \usepackage[utf8]{inputenc}
  \usepackage{textcomp} % provide euro and other symbols
\else % if luatex or xetex
  \usepackage{unicode-math} % this also loads fontspec
  \defaultfontfeatures{Scale=MatchLowercase}
  \defaultfontfeatures[\rmfamily]{Ligatures=TeX,Scale=1}
\fi
\usepackage{lmodern}
\ifPDFTeX\else
  % xetex/luatex font selection
  \setmainfont[]{Latin Modern Roman}
  \setsansfont[]{DejaVu Sans}
  \setmonofont[]{DejaVu Sans Mono}
  \setmathfont[]{Latin Modern Math}
\fi
\IfFileExists{upquote.sty}{\usepackage{upquote}}{}
\IfFileExists{microtype.sty}{% use microtype if available
  \usepackage[]{microtype}
  \UseMicrotypeSet[protrusion]{basicmath} % disable protrusion for tt fonts
}{}
\makeatletter
\@ifundefined{KOMAClassName}{% if non-KOMA class
  \IfFileExists{parskip.sty}{%
    \usepackage{parskip}
  }{% else
    \setlength{\parindent}{0pt}
    \setlength{\parskip}{6pt plus 2pt minus 1pt}}
}{% if KOMA class
  \KOMAoptions{parskip=half}}
\makeatother
\usepackage{xcolor}
\usepackage[letterpaper,margin=0.84in]{geometry}
\usepackage{longtable,booktabs,array}
\usepackage{calc} % for calculating minipage widths
\usepackage{etoolbox}
\makeatletter
\patchcmd\longtable{\par}{\if@noskipsec\mbox{}\fi\par}{}{}
\makeatother
\IfFileExists{footnotehyper.sty}{\usepackage{footnotehyper}}{\usepackage{footnote}}
\makesavenoteenv{longtable}
\setlength{\emergencystretch}{3em} % prevent overfull lines
\providecommand{\tightlist}{%
  \setlength{\itemsep}{0pt}\setlength{\parskip}{0pt}}
\setcounter{secnumdepth}{5}
\ifLuaTeX
\usepackage[bidi=basic]{babel}
\else
\usepackage[bidi=default]{babel}
\fi
\babelprovide[main,import]{american}
\ifPDFTeX
\else
\babelfont{rm}[]{Latin Modern Roman}
\fi
\let\LanguageShortHands\languageshorthands
\def\languageshorthands#1{}

\usepackage{fancyhdr}

\usepackage{microtype}
\usepackage{longtable,booktabs,array}
\definecolor{Ink}{HTML}{173345}
\definecolor{Accent}{HTML}{176B77}
\pagestyle{fancy}
\fancyhf{}
\setlength{\headheight}{15pt}
\setlength{\emergencystretch}{3em}
\setlength{\parskip}{5pt}
\setlength{\parindent}{0pt}
\allowdisplaybreaks[2]
\ifLuaTeX
  \usepackage{selnolig}  % disable illegal ligatures
\fi
\IfFileExists{bookmark.sty}{\usepackage{bookmark}}{\usepackage{hyperref}}
\IfFileExists{xurl.sty}{\usepackage{xurl}}{} % add URL line breaks if available
\urlstyle{same}



\usepackage{etoolbox}
\makeatletter
\providecommand{\subtitle}[1]{% add subtitle to \maketitle
  \apptocmd{\@title}{\par {\large #1 \par}}{}{}
}
\makeatother
\title{Wilson current concentration and conditional Gaussian fluctuations}\author{}\date{}
\begin{document}\maketitle
\hypertarget{part-a.-exact-wilson-current-concentration-and-the-extensive-bridge}{%
\section{Part A. Exact Wilson current concentration and the extensive
bridge}\label{part-a.-exact-wilson-current-concentration-and-the-extensive-bridge}}

\hypertarget{source-and-scope}{%
\subsection{Source and scope}\label{source-and-scope}}

Use the original compact product Haar Wilson law

\[
 d\mu_\beta=Z_\beta^{-1}e^{-\beta S}\,d\lambda,
 \qquad S=\sum_{p\in\mathcal P}(N-\operatorname{Re}\operatorname{Tr}U_p),
 \qquad \beta\ge1.
\]

The isotropic periodic lattice is four dimensional; all elementary
plaquettes have four distinct links. Use the product bi-invariant
Riemannian metric whose orthonormal invariant derivatives are the
original source derivatives. Write \(P=\#\mathcal P\), \(F=\nabla S\),
and let \(\Delta\) have the convention
\(\Delta\operatorname{Re}\operatorname{Tr}U=-C_F\operatorname{Re}\operatorname{Tr}U\).
The rank \(N\), metric normalization, and lattice dimension remain
fixed. Define

\[
 T=\Delta S-\beta|\nabla S|^2=3\tau_1^0.
\]

This is exactly the current in [RS], equations (3)--(4), before its physical factor
\(c_\beta^{-1}a^{-2}\). The ordinary source is still
\(R_1^0=\beta D-\ell/3\), and \(h_1^0=R_1^0+\tau_1^0\).

\hypertarget{main-weighted-theorem-conditional-scores-and-finite-colour-holder}{%
\subsection{Main weighted theorem: conditional scores and finite-colour
Holder}\label{main-weighted-theorem-conditional-scores-and-finite-colour-holder}}

For deterministic real weights \(f_e\), constant across the original
colour derivatives at each link, define

\[
 T_e=\Delta_eS-\beta|\nabla_eS|^2,\quad
 T_f=\sum_e f_eT_e,\quad M_f=\sum_e f_e^2,\quad b_f=\max_e|f_e|.
\]

There are constants \(v_0,c_0\), depending only on the fixed group and
metric, and an absolute lattice colour count \(K\le19\), such that

\[
 \boxed{\quad
 \log\mathbb E e^{tT_f}
 \le \frac{K v_0 M_f t^2}{2(1-Kc_0b_f|t|)},
 \qquad Kc_0b_f|t|<1.
 \quad} \tag{W1}
\]

No good-bath condition, gauge fixing, spectral gap, large-volume
approximation, or local Wilson energy-moment premise is used.

\hypertarget{proof-at-one-arbitrary-conditioned-link}{%
\subsubsection{Proof at one arbitrary conditioned
link}\label{proof-at-one-arbitrary-conditioned-link}}

Fix \textbf{every} other original link. The conditional law of \(U_e\)
is the genuine one-link Wilson density with its actual normalizer. Let
\(S_e\) be its incident action, up to the unchanged exterior constant.
Set \(F_e=\nabla_eS\) and \(\ell_e=\Delta_eS\). The complete compact
identity is

\[
 \mathbb E(T_e\mid\text{bath})=0,
 \quad \nabla_e\ell_e=-C_FF_e,
 \quad \ell_e\le A_*:=6C_FN,
 \quad \|\operatorname{Hess}_{ee}S\|_{op}\le H_*.
\]

The last two bounds hold uniformly for every real bath, including
cancelled staples. With \(A_e=\beta|F_e|^2=\ell_e-T_e\),

\[
 |\langle F_e,\nabla_eT_e\rangle|
 \le (C_F/\beta+2H_*)A_e\le c_0A_e,
 \quad A_e\le A_*-T_e.
\]

Apply the exponential integration-by-parts proof of the uniform-current
argument below, conditionally, with \(A_0=A_*\), \(c=c_0\). It yields

\[
 \log\mathbb E(e^{sT_e}\mid\text{bath})
 \le A_*\left[s-c_0^{-1}\log(1+c_0s)\right]
 \le\frac{v_0s^2}{2(1-c_0|s|)},\qquad v_0=A_*c_0. \tag{W2}
\]

\hypertarget{product-within-each-colour-then-holder-across-colours}{%
\subsubsection{Product within each colour, then Holder across
colours}\label{product-within-each-colour-then-holder-across-colours}}

Join two original links if a plaquette contains both. Each link is in
six plaquettes and sees at most eighteen other links, so a deterministic
proper colouring uses \(K\le19\) colours \(\mathcal C_j\). No plaquette
reads two links of one colour. Conditional on the complement of
\(\mathcal C_j\), the original Wilson density is exactly a product and
each \(T_e\), \(e\in\mathcal C_j\), reads only its own free link and
that bath. Therefore (W2) gives

\[
 \mathbb E\exp\left(t\sum_{e\in\mathcal C_j}f_eT_e\right)
 \le\exp\left\{\frac{v_0t^2\sum_{e\in\mathcal C_j}f_e^2}{2(1-c_0b_f|t|)}\right\}. \tag{W3}
\]

The bound is independent of the conditioned bath, so the outer
expectation costs nothing. Generalized Holder, with all exponents \(K\),
gives

\[
 \mathbb E e^{t\sum_jT_{f,j}}
 \le\prod_j[\mathbb E e^{KtT_{f,j}}]^{1/K},
\]

and substitution of (W3) proves (W1). The colours themselves need not be
independent.

\hypertarget{complete-rare-history-consequences}{%
\subsubsection{Complete rare-history
consequences}\label{complete-rare-history-consequences}}

For \(0\le H\le1\), \(\mathbb EH=p\),

\[
 \boxed{\quad
 \mathbb E[H T_f^2]
 \le Cp\{M_f\log(e/p)+b_f^2\log^2(e/p)\}.
 \quad} \tag{W4}
\]

Furthermore

\[
 \|T_f\|_q\le C\{\sqrt{qM_f}+qb_f\},\qquad q\ge2. \tag{W5}
\]

For every actual completed positive source-free observation, conditional
Jensen preserves (W1) for its conditional mean. The
single-ordinary-label rare refinement and its whole predictably owned
later martingale lineage have an L2 norm squared bounded by the right
side of (W4), with a fixed enlarged constant. A final orthogonal
projection contracts the L2 norm of that already rare refinement or
lineage. Such a projection does not automatically preserve the
exponential bound or a newly rare-weighted tail for an arbitrary
projected variable. The proof is exactly the rare-history argument below
with \(P\) replaced by \(M_f\) and the unit Bernstein scale replaced by
\(b_f\). Multiple common labels keep their separate
between-common-history term. A source-dependent or outcome-selected
weight is outside the deterministic-weight theorem unless separately
justified.

The original current is \(\tau_f=T_f/3\); its physical factor is still
\(c_\beta^{-1}a^{-2}\). A coefficient deterministic on the complete
probability space multiplies the bound by its squared modulus. The
complete parent mass multiplies it once, with the rare factor already
included. A coefficient depending on history needs its own envelope or
must stay inside the actual weighted sum. No independently extracted
vacuum scalar has been normalized away.

This is an \(\ell^2\) spatial bound for the \textbf{complete current}
through arbitrary actual positive history observations. It does not
apply to the ordinary \(\beta D+kS-C_0\): its directional-square terms
are not conditional one-link scores and have nonzero bath means.

\hypertarget{admitted-live-coefficient-extension}{%
\subsubsection{Admitted live-coefficient
extension}\label{admitted-live-coefficient-extension}}

The deterministic-weight proof extends to certain live coefficients. For
each colour \(\mathcal C_j\), require every \(f_e\),
\(e\in\mathcal C_j\), to be measurable with respect to the original
links outside that \textbf{whole colour}. Suppose, almost surely,
\(\sum_e f_e^2\le M\) and \(\max_e|f_e|\le b\), with deterministic
\(M,b\). Conditional product still applies with these now-fixed
coefficients, and the same Holder proof yields (W1) with \(M_f,b_f\)
replaced by \(M,b\). The weighted sum is centered since each conditional
score is centered. A bounded finite determining stencil for the
coefficients can be accommodated by a finer fixed colouring; its actual
degree must be counted. An inherited coefficient with growing/nonlocal
readset does not automatically have such a fixed colouring or a
deterministic quadratic budget.

\hypertarget{uniform-global-sub-gamma-theorem-for-the-current}{%
\subsection{Uniform global sub-gamma theorem for the
current}\label{uniform-global-sub-gamma-theorem-for-the-current}}

There are constants \(C_\Delta,C_H,c\), depending only on the fixed
group and incidence, such that

\[
 \Delta S=4C_F(NP-S),\quad \nabla\Delta S=-4C_F\nabla S,
 \quad \|\operatorname{Hess}S\|_{\mathrm{op}}\le C_H,
 \quad A_0:=4C_FNP=C_\Delta P.
\]

The Hessian bound has \textbf{no volume factor}. Each single-plaquette
Hessian is bounded on a fixed compact four-link space; summing its
quadratic-form bound counts each link only six times. This argument
concerns the Riemannian Hessian, including its connection terms, not an
unqualified matrix of noncommuting invariant derivatives.

Set \(A=\beta|F|^2\) and choose

\[
 c\ge 4C_F+2C_H>0,\qquad v=cA_0.
\]

Then \(T\le A_0\), \(\mathbb E T=0\), and for every real \(t>-1/c\),

\[
 \boxed{\quad
 \log\mathbb E e^{tT}
 \le A_0\left[t-\frac1c\log(1+ct)\right].
 \quad} \tag{1}
\]

Consequently, for \(|t|<1/c\),

\[
 \log\mathbb E e^{tT}\le\frac{vt^2}{2(1-c|t|)}, \tag{2}
\]

and for every \(x\ge0\),

\[
 \Pr(T\ge\sqrt{2vx})\le e^{-x},\qquad
 \Pr(T\le-\sqrt{2vx}-cx)\le e^{-x}. \tag{3}
\]

In particular \(\operatorname{Var}T\le v=CP\), and for \(q\ge2\),

\[
 \|T\|_q\le C\left(\sqrt{Pq}+q\right). \tag{4}
\]

\hypertarget{proof}{%
\subsubsection{Proof}\label{proof}}

For the original compact boundaryless product group, weighted
integration by parts is

\[
 \mathbb E[T f]=-\mathbb E\langle F,\nabla f\rangle.
\]

All finite-cutoff functions are smooth and bounded, so exponential
testing and differentiation require no limiting argument. Moreover

\[
 \nabla T=-4C_FF-2\beta(\operatorname{Hess}S)F,
 \qquad |\langle F,\nabla T\rangle|\le cA.
\]

Let \(M(t)=\mathbb Ee^{tT}\) and \(u(t)=(\log M)'(t)\). The exact
identity gives

\[
 M'(t)=-t\,\mathbb E[\langle F,\nabla T\rangle e^{tT}],
 \qquad |u(t)|\le c|t|\,\mathbb E_tA\le c|t|(A_0-u(t)). \tag{5}
\]

For \(t\ge0\), (5) implies \(u(t)\le A_0ct/(1+ct)\). For \(-1/c<t<0\),
it implies \(u(t)\ge A_0ct/(1+ct)\). Integrating from 0 to \(t\), with
the reversed integration direction when \(t<0\), yields (1). The
elementary bounds on \(x-\log(1+x)\) give (2), while the positive
half-axis even gives \(\log M(t)\le vt^2/2\) for all \(t\ge0\). Chernoff
optimization proves (3), and tail integration proves (4).

The argument uses neither a spectral gap nor convexity of the Wilson
action. In particular it permits all global modes and history means.

\hypertarget{an-extensive-rare-history-current-budget}{%
\subsection{An extensive rare-history current
budget}\label{an-extensive-rare-history-current-budget}}

For any random weight \(0\le H\le1\) with \(\mathbb EH=p\), (3) and
layer cake imply

\[
 \boxed{\quad
 \mathbb E[H T^2]
 \le C p\left[P\log(e/p)+\log^2(e/p)\right].
 \quad} \tag{6}
\]

The expression is zero at \(p=0\). One way to prove (6) is to use
\(\Pr(|T|>C[\sqrt{Px}+x])\le2e^{-x}\), split the layer-cake integral at
\(x=\log(2/p)\), and integrate the remaining exponential tail. This
retains one factor \(p\) and the literal original root count \(P\).

Any complete positive source-free augmentation retaining the Wilson
marginal has the same result. For an arbitrary actual observation \(Y\),
conditional Jensen preserves the exponential bound for
\(\mathbb E(T\mid Y)\). Centered differences \(T-\mathbb E(T\mid W)\)
have the same type of bound with fixed enlarged constants: use the
triangle inequality and the two tail bounds, or Hölder on the two
exponential moments. Therefore the rare-refinement proof in [RS],
Theorem 2.2, yields

\[
 \|\mathbb E(T\mid W,\alpha)-\mathbb E(T\mid W)\|_2^2
 \le Cp\left[P\log(e/p)+\log^2(e/p)\right] \tag{7}
\]

when the refinement has one ordinary label and rare complementary labels
of total mass \(p\). If there are several common labels, their
between-common-history variance remains, exactly as in [RS].

For a rare birth event \(E\in\mathcal F_b\),
\(M_t=\mathbb E(T\mid\mathcal F_t)\), and predictable ownership masks
\(a_t\in\{0,1\}\), \(a_t\le1_E\), martingale orthogonality gives

\[
 \left\|\sum_{t>b}a_t(M_t-M_{t-1})\right\|_2^2
 \le\mathbb E[1_ET^2]
 \le Cp\left[P\log(e/p)+\log^2(e/p)\right]. \tag{8}
\]

The birth increment adds once, with (7). Every subsequent actual
conditional expectation and every orthogonal output projection,
including zero-energy or low-spectral projections, contracts this bound.
No inverse spectral threshold appears. The corresponding absolute
integrated two-current connected-increment sum follows by
Cauchy--Schwarz in probability and time, retaining deterministic
physical coefficients; random history coefficients require their own
identified envelopes.

The stronger weighted result in the main weighted theorem improves the
generic \((\sum_e|f_e|)^2\) estimate for the \textbf{complete
total-current row} to its extensive \(M_f\) scale. It does not make the
same improvement for \(R_1^0\), nor for an arbitrarily weighted
directional source. Separate signed localized representatives still need
identification with the complete pushed current; all effective-score,
cutoff, live-reference and original denominator terms needed for that
identification remain.

\hypertarget{exact-covariance-locality-for-individual-current-roots}{%
\subsection{Exact covariance locality for individual current
roots}\label{exact-covariance-locality-for-individual-current-roots}}

Let

\[
 T_e=\Delta_eS-\beta|\nabla_eS|^2,
 \quad T_f=\sum_e f_eT_e,
 \quad M_f=\sum_e|f_e|^2,
\]

with field-independent real weights constant in the color index. Then

\[
 \mathbb ET_e=0,
 \qquad \mathbb E T_eT_d=0
 \quad\text{if no original plaquette contains both }e,d, \tag{9}
\]

and

\[
 \mathbb ET_e^2\le C,
 \qquad \|T_f\|_2^2\le C M_f. \tag{10}
\]

Proof of (9):
\(\mathbb E(T_eg)=-\mathbb E\langle\nabla_eS,\nabla_eg\rangle\), while
\(T_d\) reads only the original plaquettes incident to \(d\). If none
contains \(e\), its \(e\)-derivative vanishes identically. This is exact
covariance locality of current roots, not independence or mixing of the
Wilson law.

For (10), the one-link identity is \(\nabla_e\Delta_eS=-C_F\nabla_eS\).
Consequently

\[
 \mathbb ET_e^2
 =C_F\mathbb E|\nabla_eS|^2
 +2\beta\mathbb E\operatorname{Hess}_{ee}S(\nabla_eS,\nabla_eS).
\]

The absolute Hessian bound is uniform, and
\(\beta\mathbb E|\nabla_eS|^2=\mathbb E\Delta_eS\le C\). Thus
\(\mathbb ET_e^2\le C\). The bounded number of plaquette-neighbor roots
and \(2|f_ef_d|\le f_e^2+f_d^2\) prove (10).

Equivalently, the current covariance matrix on the literal original
source space satisfies \(\operatorname{Cov}(T_e,T_d)\preceq CI\). Every
actual completed positive observation \(Y\) preserves this operator
bound for
\(\operatorname{Cov}(\mathbb E(T_e\mid Y),\mathbb E(T_d\mid Y))\), by
conditional-expectation contraction. Thus the extensive arbitrary-weight
current bound propagates through every such observation including
history means.

\hypertarget{independent-proof-of-the-extensive-bridge}{%
\subsection{A direct proof of the extensive
bridge}\label{independent-proof-of-the-extensive-bridge}}

This subsection uses the local original Wilson action-moment premise
stated in [RS], equation (2), which is the local-moment theorem of
\texttt{local\_contacts/odd\_torus\_transfer\_moment.tex}: on every fixed literal plaquette record \(Y\),
\(\|\beta\sum_{p\in Y}s_p\|_q\le Cq\#Y\). It does not reprove the
reflection-positivity input.

Write \(h_f=\sum_b f_bh_b\), where
\(h_b=\beta\sum_A[\sum_\mu F_{bA;\mu}^2-F_{bA}^2/3]\), including the
fixed original pointing/averaging incidence. Compact Wilson force bounds
imply

\[
 |h_b|\le C\beta\sum_{p\in Y_b}s_p,
 \quad \|h_b\|_2\le C, \tag{11}
\]

where \(Y_b\) is a fixed-size complete determining record.

Choose any deterministic family \(\mathcal F\) of free links separated
sufficiently far that no complete \(h_b\) stencil reads two free links,
and condition on all other links, denoting that observation by \(X\).
The Wilson conditional law is exactly a product on \(\mathcal F\). The
source splits exactly as

\[
 h_f=h_{f,0}(X)+\sum_{e\in\mathcal F}H_{f,e}(U_e;X),
\]

where each \(H_{f,e}\) is the sum of \textbf{all} rooted terms whose
original stencil reads \(e\). Conditional independence gives

\[
 \mathbb E\operatorname{Var}(h_f\mid X)
 =\sum_e\mathbb E\operatorname{Var}(H_{f,e}\mid X)
 \le\sum_e\mathbb EH_{f,e}^2
 \le C M_f. \tag{12}
\]

The last inequality uses finite-stencil Cauchy--Schwarz and bounded
original incidence. No native target dimension substitutes for this
literal count.

Let \(R_f=h_f-\tau_f\), with \(\tau_f=T_f/3\), and
\(R_{X,f}=\mathbb E(R_f\mid X)\). Define

\[
 \epsilon_h=(I-P_X)h_f+P_X\tau_f,
 \qquad \epsilon_R=\epsilon_h-\tau_f=(I-P_X)R_f.
\]

Then exactly

\[
 h_f=R_{X,f}+\epsilon_h,
 \qquad R_f=R_{X,f}+\epsilon_R,
 \qquad \mathbb E\epsilon_h=\mathbb E\epsilon_R=0,
 \qquad \|\epsilon_h\|_2^2+\|\epsilon_R\|_2^2\le C M_f. \tag{13}
\]

Here the two summands in \(\epsilon_h\) are orthogonal; (10) and (12)
supply the bound. For any later actual positive observation \(Y\), even
if it reads original eliminated variables, and every orthogonal output
projection \(Q_Y\),

\[
 \|Q_Y\mathbb E(\epsilon_h\mid Y)\|_2^2+
 \|Q_Y\mathbb E(\epsilon_R\mid Y)\|_2^2\le C M_f. \tag{14}
\]

This is proved here from the inputs listed in the source locators. It requires no screening condition between \(X\) and \(Y\).
Its error remains an \(L^2\)-norm error; it cannot be converted into a
bounded difference of squared retained norms without controlling the
ordinary retained energy.

\hypertarget{why-this-does-not-solve-the-ordinary-problem}{%
\subsection{Why this does not solve the ordinary
problem}\label{why-this-does-not-solve-the-ordinary-problem}}

For the uniform test,

\[
 R=\beta D+kS-C_0,\qquad k=4C_F/3.
\]

The new argument closes for \(T\) because it is an exact
\(\mu_\beta\)-divergence and because the Wilson Laplacian closes on
\(S\): \(\nabla\Delta S=-4C_F\nabla S\). Replacing \(\beta|\nabla S|^2\)
by \(\beta D\) breaks this identity. Separate directional forces are not
the score of the complete Wilson density. The rooted ordinary
conditional mean in equation (23) of [RS] is the nonzero transverse-staple
quantity, so it cannot be erased by an invocation of the total-current
Ward identity.

The missing estimate is still an extensive bound for the actual retained
ordinary part (or its specified low-spectrum projection).
Between-history ordinary means remain. The current sub-gamma theorem
therefore supplies a real improvement to the complete error/current
sector and its rare-history propagation, while leaving the exact
ordinary obstruction exposed.

There is an actual SU(2) star demonstrating why the conditional-score
theorem cannot simply be applied to the ordinary mean. In the quaternion
convention of equation (23) of [RS], take the realizable directional
staple pairs \(h_1=2e_1\), \(h_2=-2e_1\), \(h_3=0\). The six individual
staples may be \(e_1,e_1,-e_1,-e_1,e_2,-e_2\), realized by the original
neighboring links as described in [RS], Section 5.4. Their total staple
is zero, so the conditional central-link action is constant and
\(T_e\equiv0\). Nevertheless the ordinary rooted mean is

\[
 \Phi_e=\frac34c_X^2\beta\sum_{\mu}|h_\mu|^2=6c_X^2\beta.
\]

This is a counterexample to a bathwise identification or uniform
domination by the central conditional score, not a counterexample to
extensive global susceptibility. Such stars have exponentially small
actual Wilson mass. Their aggregation, and especially the nonzero
ordinary good-star transverse means, remain additional questions.

\newpage

\hypertarget{part-b.-exact-flat-current-transfer-anchor-moments-and-its-locality-limit}{%
\section{Part B. Exact flat current transfer, anchor moments, and its
locality
limit}\label{part-b.-exact-flat-current-transfer-anchor-moments-and-its-locality-limit}}

\hypertarget{what-this-adds-to-the-uniform-averaging-rows}{%
\subsection{What this adds to the uniform averaging
rows}\label{what-this-adds-to-the-uniform-averaging-rows}}

The note [UA] proves that on a complete uniform old plaquette hierarchy
the literal flat differential has

\[
(Q_2F)_{\mu\nu}=Q_0L_\mu L_\nu F_{\mu\nu},
\qquad D_{r,\rho}^{+}Q_2F=Q_2L_\rho D_{a,\rho}^{+}F.
\]

As [UA] notes, a bound on the already summed fine divergence
alone cannot simply be applied to the separately filtered directional
terms. The cancellation below supplies a uniform operator estimate for
that summed current on the complete periodic uniform flat channel. It
also retains an arbitrary flat Bianchi defect. It is not an estimate for
the boundary/junction remainder (equation (2.12) of [UA]), the nonlinear averaging
remainders (equations (4.1)--(4.2) of [UA]), or the old constrained minimum.

The exact raw averaging map, the input/output meshes and physical
counting measures are unchanged. No reassociation of nonlinear means is
used.

\hypertarget{definitions-and-theorem}{%
\subsection{Definitions and theorem}\label{definitions-and-theorem}}

Work first on a periodic d-dimensional cubical lattice of spacing a and
side T, where r=Ba, B is an odd positive integer, and r divides T. All
forms below have values in a fixed finite-dimensional real/complex
normed vector space; Lie algebra components are handled componentwise.
Cochain components are represented at their base vertex. Antisymmetry
fixes other orientations.

Write tau\_\{s mu\} for translation by s e\_mu, and set

\[
D_\mu^+=(\tau_{a\mu}-I)/a,\quad
D_\mu^-=(I-\tau_{-a\mu})/a,\quad
H=-\sum_\mu D_\mu^-D_\mu^+,
\]

\[
L_\mu=B^{-1}\sum_{j=0}^{B-1}\tau_{ja\mu},\quad
M_\mu=L_\mu^*L_\mu,\quad
D_{r,\mu}^+=L_\mu D_\mu^+,\quad
D_{r,\mu}^-=L_\mu^*D_\mu^-.
\]

Let Q\_0 be the normalized centered B\^{}d block average. In particular
its coordinate factor in direction rho is a fixed translation of L\_rho.
Let

\[
H_r=-\sum_\mu D_{r,\mu}^-D_{r,\mu}^+
    =-\sum_\mu M_\mu D_\mu^-D_\mu^+.
\]

For an arbitrary antisymmetric 2-cochain F define its \textbf{flat}
current and Bianchi defect by

\[
J_\nu=-\sum_\mu D_\mu^-F_{\mu\nu},
\]

\[
K_{\rho\mu\nu}
=D_\rho^+F_{\mu\nu}-D_\mu^+F_{\rho\nu}+D_\nu^+F_{\rho\mu}.
\]

Finally set

\[
(J_r)_\nu=-\sum_\mu D_{r,\mu}^-(Q_0L_\mu L_\nu F_{\mu\nu}),
\]

evaluated on the actual r-grid.

\textbf{Theorem.} For each 0 \textless= eta \textless{} 2 there are
translation-invariant matrix kernels T\^{}J and T\^{}K, depending on
a,r,T and the actual block filters but not F, such that

\[
J_r=T^J J+T^K K,
\]

and, with torus distance,

\[
\sup_x\sum_y (1+\operatorname{dist}_T(x,y)/r)^\eta
 (|T^J(x,y)|+|T^K(x,y)|)\le C_{d,\eta}.
\tag{MA1}
\]

The kernels are written as coefficients of a counting-measure
convolution. The corresponding physical column bound is

\[
r^d\sup_y\sum_{x\in r\mathbb Z^d/T\mathbb Z^d}
 (1+\operatorname{dist}_T(x,y)/r)^\eta
 (|T^J(x,y)|+|T^K(x,y)|)
\le C_{d,\eta}a^d.
\tag{MA2}
\]

Consequently, for every 1 \textless= p \textless= infinity,

\[
\|J_r\|_{L_r^p}\le C_d(\|J\|_{L_a^p}+\|K\|_{L_a^p}).
\tag{MA3}
\]

For closed F, K=0. In particular an actual bound sup\textbar J\textbar{}
\textless= A in this channel gives sup r\^{}3\textbar J\_r\textbar{}
\textless= C r\^{}3 A in d=4. There is no loss r\^{}\{-1\} and no need
to estimate each directional current separately. The theorem does not
infer that the actual mixed selected fields satisfy
sup\textbar J\textbar{} \textless= C g\_k q.

\hypertarget{exact-algebra-including-zero-momentum}{%
\subsection{Exact algebra, including zero
momentum}\label{exact-algebra-including-zero-momentum}}

Let G=H\^{}\{-1\} on mean-zero cochains and zero on constant cochains.
With the above cochain conventions, the flat Hodge identity gives

\[
F_{\mu\nu}=D_\mu^+GJ_\nu-D_\nu^+GJ_\mu
 -\sum_\rho D_\rho^-GK_{\rho\mu\nu}+(P_0F)_{\mu\nu}.
\tag{MA4}
\]

One may check this by the identity d delta + delta d = H componentwise.
Every derivative commutes with H. Harmonic 2-forms on the periodic
cubical lattice are precisely constant components; the coarse current
kills them. Each J component has mean zero and satisfies sum\_mu
D\_mu\^{}- J\_mu=0. Thus no zero-mode inverse is introduced.

Substituting (MA4) into the literal coarse current gives

\[
(J_r)_\nu=Q_0L_\nu\left\{
 H_r GJ_\nu
 +D_\nu^+\sum_\mu M_\mu D_\mu^-GJ_\mu
 +\sum_{\mu,\rho}M_\mu D_\mu^-D_\rho^-GK_{\rho\mu\nu}
\right\}.
\tag{MA5}
\]

This formula includes the staggered forward/backward factors exactly. In
particular the same M\_mu = L\_mu\^{}*L\_mu, not L\_mu\^{}2, occurs in
every term.

The second term in braces can replace M\_mu by M\_mu-I because delta
J=0. Antisymmetry of K in rho,mu gives

\[
\sum_{\mu,\rho}M_\mu D_\mu^-D_\rho^-K_{\rho\mu\nu}
=\frac12\sum_{\mu,\rho}(M_\mu-M_\rho)
D_\mu^-D_\rho^-K_{\rho\mu\nu}.
\tag{MA6}
\]

The identity component cancels in (MA6). This cancellation supplies two
additional low-frequency derivatives and is the reason the endpoint
L-infinity bound is possible.

\hypertarget{heat-estimates-with-all-scales-explicit}{%
\subsection{Heat estimates with all scales
explicit}\label{heat-estimates-with-all-scales-explicit}}

Let E\_t=exp(-tH), first on the infinite a-grid. For a fixed derivative
word D\^{}alpha of length m, and 0 \textless= eta \textless= 2,

\[
\sum_z(1+|z|/\sqrt t)^\eta
 |D^\alpha E_t(z)|\le C_{d,m}t^{-m/2},\qquad t\ge a^2.
\tag{MA7}
\]

Also, for 0 \textless= t \textless= r\^{}2 and a \textless= r, E\_t is a
probability kernel with

\[
\sum_z(1+|z|/r)^\eta E_t(z)\le C_d,\qquad 0\le\eta\le2.
\tag{MA8}
\]

Here is an elementary verification sufficient for this note. The
one-dimensional kernel at u=t/a\^{}2 is symmetric and unimodal. This
follows by taking the Euler-product limit of convolution powers of the
symmetric unimodal probability kernel (u/n,1-2u/n,u/n), with n large
enough, and using closure of symmetric unimodal sequences under
convolution. Its mass at zero satisfies p\_u(0) \textless=
C(1+u)\^{}(-1/2), directly from the Fourier integral of exp(-4u
sin\^{}2(theta/2)). Telescoping yields sum\_k
\textbar p\_u(k+1)-p\_u(k)\textbar=2p\_u(0). A second telescoping gives

\[
\sum_{k\in\mathbb Z}k^2|p_u(k+1)-p_u(k)|
=p_u(0)+2\mathbb E|X_u|
\le1+2\sqrt{2u}.
\]

For u \textgreater= 1 these imply the first derivative version of (MA7),
with exponent eta=2. Product structure gives the d-dimensional result.
Fixed higher derivatives follow by splitting the heat time into m equal
parts and using weighted convolution; the bounded interval a\^{}2
\textless= t \textless= m a\^{}2 follows directly from the
difference-operator norm. Smaller eta are dominated by eta=2 at the same
sqrt(t) scale. The second moment E\textbar a X\textbar\^{}2=2dt proves
(MA8). Periodizing an infinite-lattice kernel cannot increase these
weighted norms when torus distance is used, because dist\_T(z)
\textless= \textbar z+nT\textbar. Thus the bounds are uniform in finite
torus size, including times larger than T\^{}2.

Each M\_mu is a symmetric probability average supported at distances
less than r in direction mu. The exact discrete second-difference
identity expresses

\[
M_\mu-I=\sum_q c_{\mu,q}\tau_{qa\mu}D_\mu^-D_\mu^+,
\quad c_{\mu,q}\ge0,\quad\sum_qc_{\mu,q}\le C r^2,
\quad |qa|\le r.
\tag{MA9}
\]

For example tau\_\{sa\}+tau\_\{-sa\}-2I = a\^{}2
sum\_\{q=-(s-1)\}\^{}\{s-1\}(s-\textbar q\textbar) tau\_\{qa\}
D\textsuperscript{-D}+. Averaging this identity with the symmetric
triangular weights of M proves (MA9).

\hypertarget{construction-and-bound-of-the-kernels}{%
\subsection{Construction and bound of the
kernels}\label{construction-and-bound-of-the-kernels}}

Use G=int\_0\^{}infinity(E\_t-P\_0)dt only after its derivative factors;
those factors annihilate P\_0. Split the integral at r\^{}2.

For the J channel define a matrix kernel by the sum of

\[
Q_0L_\nu E_{r^2}\,\delta_{\nu\mu},
\tag{MA10a}
\]

\[
\int_0^{r^2}Q_0L_\nu
\{H_r\delta_{\nu\mu}+D_\nu^+M_\mu D_\mu^-\}E_t\,dt,
\tag{MA10b}
\]

and

\[
\int_{r^2}^\infty Q_0L_\nu
\{(H_r-H)\delta_{\nu\mu}+D_\nu^+(M_\mu-I)D_\mu^-\}E_t\,dt.
\tag{MA10c}
\]

Summed against divergence-free, mean-zero J this is exactly its part of
(MA5). The missing P\_0 in (MA10a) acts as zero on J.

For the K channel use its original integrand in (MA5) over
{[}0,r\^{}2{]} and its antisymmetrized integrand (MA6) over
{[}r\^{}2,infinity{]}. These formulas define linear kernels on arbitrary
inputs, although their agreement with (MA5) only requires the actual
compatibility conditions already proved for J and K.

Short-time J terms have kernel norm at most C/r\^{}2: H\_r is a sum of
products of two coarse differences, and

\[
L_\nu D_\nu^+M_\mu D_\mu^-
=D_{r,\nu}^+L_\mu D_{r,\mu}^-.
\]

Every coarse difference has counting kernel norm 2/r, all remaining
averaging factors have norm one and support O(r). Formula (MA8) gives
the same bound in the weighted norm of (MA1).

For short-time K terms, rho and mu are distinct whenever K\_rho mu nu is
nonzero. Use M\_mu D\_mu\^{}-=L\_mu D\_\{r,mu\}\^{}-. The rho-coordinate
factor of Q\_0 times D\_rho\^{}- is a translated coarse difference times
an averaging operator in the other directions. Hence there are again two
coarse differences, and the weighted norm is at most C/r\^{}2. This is
where a \textbf{complete} Q\_0 block, including every coordinate,
matters.

Integrating the short-time estimate over a time interval r\^{}2 gives
O(1). The heat term (MA10a) is O(1) by (MA8).

Each long-time integrand contains an M\_mu-I (or a difference of two
such terms), two further derivatives, and therefore, by (MA9), four
lattice derivatives with total coefficient O(r\^{}2) and translations
O(r). By (MA7), for t \textgreater= r\^{}2,

\[
\|\text{integrand}\|_{\ell^1((1+|z|/r)^\eta)}
\le C_{d,\eta}r^2t^{-2}(\sqrt t/r)^\eta.
\]

Consequently

\[
\int_{r^2}^\infty C_{d,\eta}r^{2-\eta}t^{-2+\eta/2}dt
\le C_{d,\eta},\qquad \eta<2.
\]

This proves (MA1). At eta=2 the displayed integral is logarithmic; no
eta=2 or exponential localization is claimed.

For (MA2), keep Q\_0 factored on the left. For any fine-grid
displacement z,

\[
\sum_{x\in r\mathbb Z^d/T\mathbb Z^d}Q_0(x-y-z)=B^{-d}.
\]

For the J channel and all long-time terms, apply this identity to the
absolute value of the kernel before Q\_0; its weighted norm has already
been bounded. Transporting the weight across Q\_0 costs a fixed factor,
because Q\_0 has support O(r). This gives (MA2) for those terms.

The short-time K channel requires a separate column calculation, because
one coordinate factor of Q\_0 was used to regularize a derivative. This
factor cannot be counted twice. Let q\_rho be the one-dimensional
normalized B-point average in that direction. Its derivative is exactly
a difference of two shifted point masses divided by r. The
one-dimensional heat kernel, sampled on a spacing-r sublattice,
satisfies for 0 \textless= t \textless= r\^{}2

\[
\sup_y\sum_{x\in r\mathbb Z}(1+|x-y|/r)^\eta E_t^{(1)}(x-y)
\le C_\eta\min\{1,a/\sqrt t\},
\tag{MA12}
\]

with the right side interpreted as C\_eta at t=0. For t \textless=
a\^{}2 the unrestricted heat-kernel moment proves this. For a\^{}2
\textless= t \textless= r\^{}2, the elementary bound

\[
E_t^{(1)}(ak)\le C(a/\sqrt t)\exp(-c|ak|/\sqrt t)
\]

proves it by a geometric series on the spacing-r grid, uniformly in the
residue y. To justify this heat bound directly, split the time in half,
use the previously proved supremum bound C a/sqrt(t), and bound the tail
of the other half by exponential Markov with parameter c/sqrt(t). Its
moment generating function is exp(2(t/a\^{}2)(cosh(ca/sqrt(t))-1)),
bounded uniformly when t \textgreater= a\^{}2. Periodization preserves
the bound with torus distance.

In the rho direction, D\_rho\^{}- q\_rho is a pair of point masses
divided by r, so (MA12) bounds its sampled heat convolution by C
r\^{}\{-1\} min(1,a/sqrt(t)). In the mu direction the other coarse
difference costs 2/r, while its Q\_0 average and those in every other
direction retain the exact B\^{}\{-1\} column factors. Here mu != rho,
since the corresponding K component is zero otherwise. Remaining L
filters shift by at most r and have mass one. Thus the entire short-time
K integrand has weighted coarse column sum at most

\[
C_{d,\eta}B^{-(d-1)}r^{-2}\min\{1,a/\sqrt t\}.
\]

Its time integral is bounded by

\[
C B^{-(d-1)}r^{-2}\left(a^2+\int_{a^2}^{r^2}a t^{-1/2}dt\right)
\le C B^{-(d-1)}a/r=C B^{-d}.
\]

This supplies the missing column count without changing the operator or
inserting an extra average. Product weights dominate the Euclidean
weight, so the coordinate calculation proves the stated d-dimensional
weighted column bound. Multiplying by r\^{}d gives (MA2). The endpoint
p=1 and p=infinity bounds now imply (MA3), by interpolation or the
weighted Schur argument.

\hypertarget{finite-parameter-jets-and-original-anchor-measures}{%
\subsection{Finite parameter jets and original anchor
measures}\label{finite-parameter-jets-and-original-anchor-measures}}

All these operators are fixed linear maps, so they commute with every
fixed parameter derivative. For an m-th response kernel keep its
original anchor sets and original counting measures, and define
A\_eta\^{}(m) as in equation (3.4) of [UA], with old physical metric dist/r.
Entrywise triangle inequality and

\[
1+d(x,z)\le(1+d(x,y))(1+d(y,z))
\]

give

\[
A_\eta^{(m)}(J_r;x)
\le\sum_y |T^J(x,y)|(1+d(x,y))^{m\eta}A_\eta^{(m)}(J;y)
 +\sum_y |T^K(x,y)|(1+d(x,y))^{m\eta}A_\eta^{(m)}(K;y).
\tag{MA11}
\]

For m eta \textless{} 2, (MA1) proves a uniform sup bound by the two
corresponding input sup norms. It does not manufacture an input anchor
sum from a directional norm. For a mixed native metric the literal
weighted kernel sum must replace (MA1); an old Euclidean range does not
imply bounded mixed-metric diameter.

\hypertarget{the-second-moment-threshold-is-a-genuine-limitation}{%
\subsection{The second-moment threshold is a genuine
limitation}\label{the-second-moment-threshold-is-a-genuine-limitation}}

For d \textgreater= 3 and B \textgreater{} 1, no convolution
representation of this current map on all compatible closed 2-forms can
have a uniformly finite absolute second spatial moment (and hence none
can have uniform exponential locality). The statement concerns
dependence on the \textbf{total fine current alone}; it does not deny
the finite stencil of the original map as a function of F.

Here is a direct multiplier proof. Choose a component nu and Fourier
momenta p with p\_nu=0 and two other momenta free. For every nonzero
such momentum, J=e\_nu j(p) is divergence free and mean zero. It is
realized by the closed cochain F=dH\^{}\{-1\}J. Since D\_nu\^{}+=0,
(MA5) reduces exactly to

\[
(J_r)_\nu(p)=q_0(p)\frac{h_r(p)}{h_a(p)}j(p),\qquad
h_s(p)=\sum_\mu\frac{4}{s^2}\sin^2(sp_\mu/2).
\tag{MA13}
\]

The L\_nu factor equals one. For physical Fourier momentum p tending to
zero,

\[
\frac{h_r(p)}{h_a(p)}
=1-\frac{r^2-a^2}{12}\,
 \frac{\sum_\mu p_\mu^4}{\sum_\mu p_\mu^2}
 +O(|p|^4).
\tag{MA14}
\]

The centered block filter has
q\_0(p)=1-(r\textsuperscript{2-a}2)\textbar p\textbar{}\textsuperscript{2/24+O(\textbar p\textbar{}}4).
Its quadratic part is analytic and cannot remove the non-polynomial
homogeneous degree-two term in (MA14). In the plane of two nonzero
momentum coordinates,
(x\textsuperscript{4+y}4)/(x\textsuperscript{2+y}2) takes the value
t\^{}2 at (t,0), (0,t), (t,t), and (t,-t). No quadratic form can agree
with those four directional limits. Thus the multiplier is not C\^{}2 at
zero.

A convolution kernel with finite absolute second moment has a C\^{}2
Fourier transform by differentiation under an absolutely summable
series, contradicting (MA13)--(MA14). Uniform finite-torus second-moment
bounds would yield such an infinite-volume kernel by a subsequence and
Fatou; the finite-torus compatible Fourier modes become dense and
identify its restriction (MA13). Taking paired opposite momenta gives
real fields if desired. Removing the zero Fourier coefficient of J does
not remove the obstruction, which uses nonzero momenta approaching zero.

Consequently eta\textless2 is not merely an artifact of the long-time
integral. This theorem supplies uniform current amplitude and finite
polynomial anchor moments within that range; it does not supply the
saved exponential anchor estimates for a general source packet.
Additional structure of the actual packet would have to establish those
estimates.

\hypertarget{remaining-admission-obligations-and-interpretation}{%
\subsection{Remaining admission obligations and
interpretation}\label{remaining-admission-obligations-and-interpretation}}

The result is a cancellation for the exact flat complete uniform old
channel. For a family of closed linearized curvatures it controls the
old total current from the fine total current with no scale-ratio loss.
For an arbitrary Lie-algebra curvature cochain it exposes the flat
Bianchi defect as a separate input. A genuinely nonabelian discrete
Bianchi identity, with actual transports and their derivatives, is not
supplied by simply writing dF=0.

To apply to the actual incoming representative of [UA], one still has to
identify the same real and simultaneous complex wall chart; bound the
fine physical current versus its flat counterpart and the Bianchi
defect; bound every nonlinear averaging/reference remainder and
boundary/junction row; prove the old minimum's Maxwell/current/anchor
responses on its saved radius; retain quadratic/cubic covariant
completion; and attach an unused reserve on the same complete selected
body. Nothing here enlarges the admitted domain or changes any
source/scalar prescription.

\newpage

\hypertarget{part-c.-actual-wilson-conditional-gaussian-fluctuations}{%
\section{Part C. Actual-Wilson conditional Gaussian
fluctuations}\label{part-c.-actual-wilson-conditional-gaussian-fluctuations}}

\hypertarget{inputs-and-exact-scope}{%
\subsection{Inputs and exact scope}\label{inputs-and-exact-scope}}

The original conventions, from [ND], equations (1.1)--(1.3), are

\[
 d\mu_{a,\beta}=Z_{a,\beta}^{-1}e^{-\beta S_0}\,d\lambda,
 \qquad O_{H,a}(f)=c_\beta^{-1}a^{-2}\beta G_{H,a}(f),
 \qquad \beta=c_\beta g^{-2}.
\]

Here \texttt{c\_beta} is the fixed convention constant, not a variable
depending on beta. The complete pointed/hypercubic test average remains
in \texttt{G\_H}. This note writes each completed rooted contribution as
a finite sum of terms \texttt{f(x\_alpha)\ h\_alpha}, with all the
literal shifts retained and with boundedly many shifts within
\texttt{C\ a} of the root. The local nonnegative quadratic directional
residual is

\[
 h_b=\sum_A\left[\sum_{\mu\ne\nu_b}(F^A_{b;\mu})^2
             -\frac13\left(\sum_{\mu\ne\nu_b}F^A_{b;\mu}\right)^2\right]\ge0.
\]

For every fixed literal plaquette record \texttt{Q}, the
local-moment theorem of \texttt{local\_contacts/odd\_torus\_transfer\_moment.tex} gives

\[
 \sup_a E\exp\left(\theta_Q\beta\sum_{p\in Q}s_p\right)<\infty
 \tag{1}
\]

on its explicitly stated physical regime

\[
 n_a\ge c_*\log(2+N\beta_a),\qquad n_a\ge3,
\]

where the original isotropic torus has \texttt{n\_a\^{}4} sites. Its partition-power input is
taken from the manuscript (Lemma 10.1); that partition theorem is not
reproved here. A
more general version retains \texttt{K\_a=1+log(2+N\ beta\_a)/n\_a} in
all local moments. The uniform-moment theorem below uses the displayed
physical regime. It does not assert a running-coupling trajectory.

The second input is the same original Wilson one-link good-bath
calculation in [ND], Sections 2.1--2.3: the complete
free-link contribution at a flat bath is

\[
 G_{e,\mathrm{flat}}(U)=12|\nabla s(U)|^2+\text{constant},
\]

and its rescaled variable \texttt{beta\ G\_e} has a fixed concentration
deficit on every sufficiently good bath. The bad-bath probability is at
most

\[
 \epsilon_a=C\frac{\log(2+\beta_a)}{\beta_a}\longrightarrow0
 \tag{2}
\]

if \texttt{beta\_a\ -\textgreater{}\ infinity}. Under (1), the same
finite-record event gives the stronger \texttt{C\ exp(-c\ beta\_a)}
bound, but that strengthening is not needed.

\hypertarget{complete-independent-stencils-and-unchanged-source-decomposition}{%
\subsection{Complete independent stencils and unchanged source
decomposition}\label{complete-independent-stencils-and-unchanged-source-decomposition}}

Fix a bounded physical cube \texttt{Omega} whose closure embeds in the
original torus, and choose one fixed link orientation. Select a
deterministic sublattice of links \texttt{F\_a} with spacing
\texttt{K\ a}, where \texttt{K} is a fixed integer larger than twice the
diameter, in fine lattice units, of every complete determining stencil.
Keep only selected links whose full determining stencil and entire
packing cell lie in \texttt{Omega}. The omitted boundary strip has
physical volume O(a). Changing this deterministic set does not change
the measure. No plaquette reads two selected links, and no completed
rooted \texttt{H} contribution reads two selected links.

Let \texttt{B\_a} be the sigma algebra generated by all original links
outside \texttt{F\_a}. Assign to \texttt{G\_\{a,e\}(f)} every complete
rooted summand that reads \texttt{e}, and assign all remaining summands
to \texttt{G\_\{a,0\}(f)}. This assignment is unique because of stencil
separation. Thus

\[
 G_{H,a}(f)=G_{a,0}(f)+\sum_{e\in F_a}G_{a,e}(f),
 \tag{3}
\]

where \texttt{G\_\{a,0\}} is \texttt{B\_a} measurable, and the actual
conditional Wilson law of the free links is a product. Every neighboring
root that varies with \texttt{e} is included. In particular this is not
the zero residual at the selected link alone.

Define

\[
 X_{a,e}(f)=\beta_aG_{a,e}(f),\quad
 Y_{a,e}(f)=X_{a,e}(f)-E(X_{a,e}(f)\mid B_a),
\]

\[
 S_a(f)=a^2\sum_{e\in F_a}Y_{a,e}(f),\qquad
 V_a(f,g)=a^4\sum_{e\in F_a}
 E(Y_{a,e}(f)Y_{a,e}(g)\mid B_a).
 \tag{4}
\]

The exact relationship to the unchanged physical source is

\[
 \boxed{S_a(f)=c_\beta a^4\big[O_{H,a}(f)-E(O_{H,a}(f)\mid B_a)\big].}
 \tag{5}
\]

There is no assertion that the bath conditional mean is a deterministic
subtraction or an admissible renormalization of the full physical field.

\hypertarget{actual-law-local-moment-bound}{%
\subsection{Actual-law local moment
bound}\label{actual-law-local-moment-bound}}

For the compact group \texttt{SU(N)} and the fixed invariant metric,

\[
 |\nabla s(U)|^2\le C_Ns(U)
\]

globally. Near the identity this follows from the positive quadratic
minimum of \texttt{s}; away from the identity it follows from
compactness and positivity of \texttt{s}. Each force reads a bounded
number of plaquettes. Consequently

\[
 |X_{a,e}(f)|\le C\|f\|_{\infty,\,\mathrm{stencil}(e)}
                  \beta_a\sum_{p\in Q_e}s_p,
 \tag{6}
\]

where \texttt{Q\_e} has uniformly bounded cardinality and is the actual
enlarged stencil. Equation (1), Holder's inequality, and conditional
Jensen yield, at every fixed \texttt{p\textgreater{}=1},

\[
 E|Y_{a,e}(f)|^p
 \le C_p\|f\|_{\infty,\,\mathrm{stencil}(e)}^p.
 \tag{7}
\]

These are averaged same-Wilson moments. They do \textbf{not} assert a
uniform conditional moment on every bath or history.

\hypertarget{conditional-gaussian-approximation-with-a-quantitative-rate}{%
\subsection{Conditional Gaussian approximation with a quantitative
rate}\label{conditional-gaussian-approximation-with-a-quantitative-rate}}

\hypertarget{lemma-1-independent-centered-triangular-arrays-retaining-random-variances}{%
\subsubsection{Lemma 1: independent centered triangular arrays,
retaining random
variances}\label{lemma-1-independent-centered-triangular-arrays-retaining-random-variances}}

Let \texttt{Y\_e} be independent conditional on a sigma algebra
\texttt{B}, with conditional means zero and finite third moments. Set
\texttt{v\_e=E(Y\_e\^{}2\textbar{}B)}. For every real \texttt{t} and
\texttt{s\textgreater{}=0},

\[
 \left|E\left(e^{it s\sum_eY_e}\mid B\right)
       -\exp\left(-\frac{t^2s^2}{2}\sum_ev_e\right)\right|
 \le C|t|^3s^3\sum_eE(|Y_e|^3\mid B).
 \tag{8}
\]

\textbf{Proof.} For each factor, Taylor's formula gives

\[
 |E(e^{iuY_e}\mid B)-(1-u^2v_e/2)|
 \le |u|^3E(|Y_e|^3\mid B)/6.
\]

For \texttt{z\textgreater{}=0},
\texttt{\textbar{}exp(-z)-1+z\textbar{}\ \textless{}=\ C\ z\^{}(3/2)}.
Also
\texttt{v\_e\^{}(3/2)\ \textless{}=\ E(\textbar{}Y\_e\textbar{}\^{}3\textbar{}B)}.
Therefore the difference between that characteristic factor and
\texttt{exp(-u\^{}2v\_e/2)} is at most
\texttt{C\ \textbar{}u\textbar{}\^{}3\ E(\textbar{}Y\_e\textbar{}\^{}3\textbar{}B)}.
Both types of factors have modulus at most one. Telescope their
products. No maximum conditional-variance hypothesis is needed. QED.

Applying (8) to (4), using \texttt{s=a\^{}2}, (7), and
\texttt{\#F\_a=O(a\^{}-4)}, proves

\[
 \boxed{
 E\left|E(e^{itS_a(f)}\mid B_a)-e^{-t^2V_a(f,f)/2}\right|
 \le C_f|t|^3a^2.
 }
 \tag{9}
\]

For a finite list of real smooth tests \texttt{f\_1,...,f\_m} and real
\texttt{t\_i}, the same result gives

\[
 E\left|E\left(e^{i\sum_it_iS_a(f_i)}\mid B_a\right)
 -\exp\left(-\tfrac12\sum_{i,j}t_it_jV_a(f_i,f_j)\right)\right|
 \le C a^2\left(\sum_i|t_i|\|f_i\|_\infty\right)^3.
 \tag{10}
\]

The constants may depend on the fixed compact common test support. In
the version with \texttt{K\_a} moments, the right side of (9)--(10) is
\texttt{C\ a\^{}2\ K\_a\^{}3}.

For every bounded \texttt{B\_a}-measurable random variable
\texttt{Z\_a}, (9) also gives

\[
 \left|E[Z_ae^{itS_a(f)}]-E[Z_ae^{-t^2V_a(f,f)/2}]\right|
 \le C_f\|Z_a\|_\infty|t|^3a^2.
 \tag{11}
\]

This is the crucial retained-bath statement. It does not replace the
variance by its expectation and does not assert mixing of the bath.

\hypertarget{conditional-cumulants}{%
\subsubsection{Conditional cumulants}\label{conditional-cumulants}}

At each fixed order \texttt{r\textgreater{}=3}, conditional independence
and the moment-cumulant formula give

\[
 E\left|\kappa_r(S_a(f_1),...,S_a(f_r)\mid B_a)\right|
 \le C_{r,f_1,...,f_r}a^{2r-4}.
 \tag{12}
\]

To check the bound, each conditional cumulant is a fixed finite sum of
products of conditional moments. Conditional Holder followed by ordinary
Holder bounds the expectation of each product by a fixed \texttt{r}th
absolute moment. Equation (7) bounds that moment. Sum the
\texttt{O(a\^{}-4)} independent contributions and retain
\texttt{a\^{}(2r)}. The first conditional cumulant is zero exactly; the
second is (4). For disjoint supports separated by more than the stencil
diameter, a mixed conditional cumulant involving both supports is zero
exactly, since no selected-link variable reads both.

\hypertarget{the-limiting-covariance-is-a-random-local-intensity-not-a-guessed-constant}{%
\subsection{The limiting covariance is a random local intensity, not a
guessed
constant}\label{the-limiting-covariance-is-a-random-local-intensity-not-a-guessed-constant}}

Write

\[
 X_{a,e}^0=\beta_aG_{a,e}(1),\qquad
 v_{a,e}=\operatorname{Var}(X_{a,e}^0\mid B_a),
 \qquad \nu_a=a^4\sum_{e\in F_a}v_{a,e}\delta_{x_e}.
 \tag{13}
\]

The notation \texttt{G\_e(1)} means the sum of the complete rooted terms
assigned to that selected link, with every test coefficient replaced by
one. It does not add roots assigned elsewhere.

The actual pointed test offsets satisfy

\[
 E|Y_{a,e}(f)-f(x_e)Y_{a,e}(1)|^2
 \le C_fa^2.
 \tag{14}
\]

Indeed each coefficient difference is at most
\texttt{C\ a\ \textbar{}\textbar{}grad\ f\textbar{}\textbar{}\_infty},
and (6)--(7) control the entire finite local row. Conditional
independence then gives

\[
 \left\|S_a(f)-a^2\sum_ef(x_e)Y_{a,e}(1)\right\|_{L^2(\mu)}
 \le C_fa,
 \tag{15}
\]

and Cauchy--Schwarz gives

\[
 E\left|V_a(f,g)-\int fg\,d\nu_a\right|\le C_{f,g}a.
 \tag{16}
\]

Thus the characteristic-function approximation with \texttt{V\_a(f,g)}
replaced by \texttt{int\ f\ g\ d\ nu\_a} has error \texttt{O(a)} for
fixed smooth tests. This error comes from the \textbf{actual} shifted
tests, not from changing their definition.

For complete packing cells \texttt{C\_\{a,e\}} of volume
\texttt{(K\ a)\^{}4}, define

\[
 w_a(x)=K^{-4}v_{a,e}\quad(x\in C_{a,e}),
\]

and put zero on the boundary strip not covered by complete cells. For
every \texttt{1\textless{}p\textless{}infinity},

\[
 E\int_\Omega w_a(x)^p\,dx\le C_{p,\Omega}.
 \tag{17}
\]

This follows from conditional Jensen and
\texttt{E\textbar{}X\_e\^{}0\textbar{}\^{}(2p)\textless{}=C\_(2p)}. The
atomic measure \texttt{nu\_a} and \texttt{w\_a\ dx} differ against a
fixed Lipschitz function by at most
\texttt{C\ a\ \textbar{}\textbar{}grad\ f\textbar{}\textbar{}\_infty\ nu\_a(Omega)}.
Their expected total masses are uniformly bounded.

Consequently their laws are tight as positive measures on compact
subsets. Every subsequential weak limit \texttt{nu} has an almost surely
locally \texttt{L\^{}p} density \texttt{w}, for every finite \texttt{p},
with the corresponding integrated moment bounds. One precise proof uses
the lower semicontinuity on measures of

\[
 \mathcal F_p(\nu)=
 \begin{cases}
 \int|w|^p dx,&\nu=w\,dx,\\
 +\infty,&\text{otherwise},
 \end{cases}
\]

which follows from its convex dual formula. Apply this to
\texttt{w\_a\ dx}, not to the atomic approximations. A diagonal
extraction over integer \texttt{p\textgreater{}=2} gives all finite
\texttt{p}. Merely bounding \texttt{E\ nu\_a} would not suffice to rule
out a singular random intensity.

\hypertarget{nondegenerate-intensity}{%
\subsubsection{Nondegenerate intensity}\label{nondegenerate-intensity}}

The good-bath concentration deficit of [ND] has constants
\texttt{h\_0,p\_0\textgreater{}0} such that every interval of length
\texttt{h\_0} carries conditional probability at most \texttt{1-p\_0}
for \texttt{X\_e\^{}0}. Taking the interval centered at its conditional
mean yields

\[
 v_{a,e}\ge v_*:=p_0h_0^2/4
 \tag{18}
\]

on every good bath. The expected physical volume of bad packing cells is
\texttt{O(epsilon\_a)}. The uncovered physical boundary strip has volume
\texttt{O(a)}. The empirical selected-link measure
\texttt{a\^{}4\ sum\_e\ delta\_(x\_e)} tends to \texttt{K\^{}-4\ dx} in
the interior. Therefore every limiting density obeys

\[
 \boxed{w(x)\ge K^{-4}v_*>0\quad\text{for almost every }x\in\Omega,\text{ almost surely}.}
 \tag{19}
\]

For example, to justify passage to the limit, test against a countable
dense set of nonnegative smooth functions compactly supported in
\texttt{Omega}; the negative defect from bad cells tends to zero in
\texttt{L\^{}1}. Equation (19) follows by density. No independence of
the good-bath events is used.

\hypertarget{actual-source-subsequential-limit-and-what-it-does-not-prove}{%
\subsection{Actual-source subsequential limit and what it does not
prove}\label{actual-source-subsequential-limit-and-what-it-does-not-prove}}

For every fixed \texttt{s\textgreater{}2}, the conditional-centering
identity and (7) give a uniform second-moment bound for \texttt{S\_a} in
\texttt{H\^{}\{-s\}} on a compact interior cube or torus: test against a
bounded Fourier/eigenfunction basis and sum
\texttt{(1+\textbar{}k\textbar{}\^{}2)\^{}(-s)}. Taking any stronger
negative index gives compact containment. Hence the fluctuation fields
have distributional subsequential limits. Equivalently, one can first
extract finite-dimensional laws on a countable dense test family and use
the same moment bound for the random-distribution realization.

Along a joint subsequential limit of the fluctuation fields, intensity
measures, and retained bath variables, where \texttt{Z\_a} is a
bath-measurable random element in a fixed Polish space and the joint
laws are tight, (10)--(11) and (16) imply the joint characteristic
functional

\[
 E\left[e^{i\sum_it_iS(f_i)}\mid\nu,Z\right]
 =\exp\left(-\frac12\int\left(\sum_it_if_i(x)\right)^2w(x)\,dx\right).
 \tag{20}
\]

Here \texttt{Z} stands for retained bath observables in the chosen joint
limit. Thus the limit is \textbf{conditionally Gaussian white noise with
a random intensity density}. Conditional on the retained bath,
disjoint-support tests are independent. They need not be independent
without conditioning. The distribution can be realized as
\texttt{sqrt(w)\ W} with an auxiliary white noise independent of
\texttt{(w,Z)}; this is a representation of the proved limit kernel, not
a claim of independence at finite cutoff.

The full original centered field still satisfies the exact decomposition

\[
 a^4[O_{H,a}(f)-EO_{H,a}(f)]
 =c_\beta^{-1}S_a(f)
  +a^4[E(O_{H,a}(f)\mid B_a)-EO_{H,a}(f)].
 \tag{21}
\]

This note gives no tightness, susceptibility estimate, Gaussian limit,
or physical identification for the second term. Nor does it identify
\texttt{w} with a constant. Those are precisely spatial/low-spectral
questions, rather than consequences of one-link conditional
independence. The random local intensity is a retained bath response; it
cannot be replaced by a field-independent scalar contact counterterm
without a separate argument.

The theorem proves an extensive actual conditional noise energy:

\[
 E\operatorname{Var}(\beta G_{H,a}(f)\mid B_a)=O(a^{-4}),
\]

with a matching lower bound for plateau tests, and identifies its
conditional higher cumulants at their exact scales. It does not bound
the between-bath variance. Consequently it does not close all-history
MT36, AE8/LM47, S94, or the common interacting continuum hierarchy. It
should be used to isolate a genuine ultraviolet contact component while
continuing to study the unchanged retained ordinary mean.

\hypertarget{flat-bath-normalization-check}{%
\subsection{Flat-bath normalization
check}\label{flat-bath-normalization-check}}

Let the metric convention give
\texttt{s(exp(x))=(kappa/2)\textbar{}x\textbar{}\^{}2+O(\textbar{}x\textbar{}\^{}4)}
in an orthonormal Lie basis, and put \texttt{d=dim\ SU(N)}. In the
exactly flat one-link bath, the conditional action is \texttt{6\ s(U)},
so \texttt{z=sqrt(beta)\ x} tends to a centered Gaussian with covariance
\texttt{(6\ kappa)\^{}-1\ I}. Therefore

\[
 \beta G_{e,\mathrm{flat}}(U)-\text{constant}
 \Longrightarrow 2\kappa\,\chi_d^2,
\]

\[
 \lim_{\beta\to\infty}v_{e,\mathrm{flat}}=8\kappa^2d.
 \tag{22}
\]

This checks the positive variance and normalization. It is a single
exactly conditioned flat bath, not a claim that the actual bath variance
converges to this constant. A nonflat bath can contribute a rescaled
linear term in the complete local source, so retaining the random
intensity in (20) is necessary until that response is controlled.

\hypertarget{stronger-version-joint-physical-and-ordinary-noises-from-two-link-blocks}{%
\subsection{Stronger version: joint physical and ordinary noises from
two-link
blocks}\label{stronger-version-joint-physical-and-ordinary-noises-from-two-link-blocks}}

The complete ordinary row can also be retained in the same theorem.
Define its dimensionless source by

\[
 h_a(f)=\beta_aG_{H,a}(f),\qquad
 R_a(f)=\beta_aD_a[f]-\ell_a[f]/3.
 \tag{23}
\]

Thus \texttt{O\_H=c\_beta\^{}-1\ a\^{}-2\ h} and the literal specified
ordinary partner is \texttt{r\_H=c\_beta\^{}-1\ a\^{}-2\ R}. The
identity and field-dependent Casimir terms are both kept.

Replace each selected one-link site by the exact pair from the companion
two-link theorem (source locator 4):

\[
 b_-=(x-e_\mu,\nu),\qquad b_+=(x+e_\mu,\nu),\qquad\mu\ne\nu.
\]

No plaquette reads both links. Their stars meet only at the middle bond
\texttt{(x,nu)}. Separate the pair centers by a fixed large lattice
spacing so every complete rooted source summand reads at most one pair.
Conditional on all other links, the free links are independent, hence
the two-link blocks are independent. For each block, retain the
\textbf{entire} variable contributions to both (23), and conditionally
center each. Call the resulting vector
\texttt{Y\_\{a,e\}=(Y\^{}h\_\{a,e\},Y\^{}R\_\{a,e\})}.

The scalar lemma applies to every real linear combination of these
vectors. Equations (9)--(17), with a symmetric positive 2 by 2
covariance matrix in place of the scalar variance, follow verbatim.
Local moments for the \texttt{R} row follow from the force-square energy
bound and the exact local Casimir formula

\[
 \ell[f]=C_F\sum_pf_p(N-s_p),
\]

with its original finite incidences. The local Casimir row is bounded
even before conditional centering; its constant drops from the variance
only because variance is invariant under this \textbf{explicit}
centering, not because it has been removed from (23).

The useful improvement is a uniformly \textbf{positive definite}
limiting covariance matrix. This requires a two-source variance lemma,
not a scalar variance lower bound applied separately to the two fields.

\hypertarget{lemma-2-variance-from-a-finite-family-of-second-derivative-tests}{%
\subsubsection{Lemma 2: variance from a finite family of second
derivative
tests}\label{lemma-2-variance-from-a-finite-family-of-second-derivative-tests}}

Suppose a probability density \texttt{p} is bounded below by
\texttt{m\_0\textgreater{}0} on a fixed coordinate box \texttt{Q}.
Suppose a smooth real function \texttt{q} has, on all of \texttt{Q}, one
of the finitely many second derivatives \texttt{L\_j\ q} of constant
sign and magnitude at least \texttt{kappa\textgreater{}0}, where
\texttt{L\_j} is either a pure or mixed coordinate second derivative.
Then

\[
 \operatorname{Var}_p(q)\ge
 \frac{m_0\kappa^2}{\|L_j\varphi\|_{L^2(Q)}^2}>0,
 \tag{24}
\]

for any fixed nonnegative \texttt{varphi\ in\ C\_c\^{}infinity(Q)} of
integral one.

\textbf{Proof.} Reverse the sign of \texttt{L\_j} if needed. Integration
by parts gives
\texttt{int\ q\ L\_j\ varphi=int\ (L\_j\ q)\ varphi\textgreater{}=kappa}.
Since \texttt{int\ L\_j\ varphi=0}, one can subtract \texttt{E\_p\ q}.
Cauchy--Schwarz and \texttt{p\textgreater{}=m\_0} imply (24). No density
derivative or bound on the linear term of \texttt{q} is used. QED.

In the companion two-link theorem's orthonormal convention
\texttt{-tr(T\^{}A\ T\^{}B)=delta\_AB}, the exact nonlinear sources have
flat quadratic jets

\[
 G_H^{(2)}(u,v)=12(|u|^2+|v|^2)+\frac43\langle u,v\rangle,
\]

\[
 D^{(2)}(u,v)=30(|u|^2+|v|^2)+2\langle u,v\rangle.
 \tag{25}
\]

For any fixed generator coordinate, the two relevant derivative rows are
therefore

\[
 \begin{pmatrix}
 \partial_{u_A}^2G_H&\partial_{u_A}^2D\\
 \partial_{u_A}\partial_{v_A}G_H&\partial_{u_A}\partial_{v_A}D
 \end{pmatrix}_{(0,0)}
 =\begin{pmatrix}24&60\\4/3&2\end{pmatrix},
 \qquad \det=-32\ne0.
 \tag{26}
\]

Use the actual conditional minima and canonical variables
\texttt{z=sqrt(beta)(u-u\_*)}, \texttt{w=sqrt(beta)(v-v\_*)}. A
canonical second derivative of \texttt{beta\ G} is exactly the
corresponding physical second derivative of \texttt{G}. For the Casimir
companion its canonical second derivative is \texttt{-1/(3\ beta)} times
the corresponding physical derivative of \texttt{ell}; those physical
derivatives are uniformly bounded on the fixed chart. Its mixed
derivative is in fact zero because no plaquette reads both free links.
Thus (26), finite-word continuity, and sufficiently small fixed
good-bath radius imply the following uniform fact: for every unit vector
\texttt{theta=(theta\_h,theta\_R)}, at least one of the two tested
derivatives of

\[
 q_\theta=\theta_h\beta G_{e,H}
          +\theta_R(\beta D_e-\ell_e/3)
\]

has constant sign and magnitude bounded below on the fixed canonical
box. One first chooses the nonzero derivative from the flat matrix,
whose smallest singular value is positive; all actual derivative errors
are then made smaller than half that margin uniformly in \texttt{theta}.
The same original conditional density minorization as the companion
two-link theorem and [ND] applies on this box. Lemma 2 proves

\[
 \operatorname{Cov}(Y_{a,e}\mid B_a)\succeq c_* I_2
 \tag{27}
\]

on every good two-link bath. This is a lower bound for all linear
combinations, not just for the two diagonal variances.

The matrix intensity measures

\[
 \boldsymbol{\nu}_a=a^4\sum_e\operatorname{Cov}(Y_{a,e}\mid B_a)\delta_{x_e}
\]

have subsequential densities \texttt{mathsf\ W(x)} with all finite local
\texttt{L\^{}p} moments and

\[
 \mathsf W(x)\succeq K^{-4}c_*I_2
 \quad\text{almost everywhere, almost surely}.
 \tag{28}
\]

To justify matrix compactness, use the trace to dominate the total
variation of every entry. Moment bounds follow from conditional Holder.
Positive semidefiniteness and (28) pass by testing a countable dense set
of rational source vectors and nonnegative smooth spatial tests. Every
joint subsequential fluctuation limit has conditional characteristic
functional

\[
 E\left[e^{iS^h(f)+iS^R(g)}\mid\mathsf W,Z\right]
 =\exp\left[-\tfrac12\int_\Omega
       (f(x),g(x))\mathsf W(x)(f(x),g(x))^T\,dx\right].
 \tag{29}
\]

Both exact physical identities remain:

\[
 S_a^h(f)=c_\beta a^4[O_{H,a}(f)-E(O_{H,a}(f)\mid B_a)],
\]

\[
 S_a^R(f)=c_\beta a^4[r_{H,a}(f)-E(r_{H,a}(f)\mid B_a)].
 \tag{30}
\]

The current companion is their exact difference. Its scalar, lower, and
source-response records are not changed. Their \textbf{between-bath
means} still require control; the two-link theorem proves no
cancellation of them.

\hypertarget{flat-two-link-covariance-check}{%
\subsubsection{Flat two-link covariance
check}\label{flat-two-link-covariance-check}}

At the flat bath \texttt{z,w} are independent Gaussian vectors with
covariance \texttt{I/6}. Put \texttt{d=N\^{}2-1},
\texttt{A=\textbar{}z\textbar{}\^{}2+\textbar{}w\textbar{}\^{}2},
\texttt{B=z\ dot\ w}. Then

\[
 \operatorname{Var}A=d/9,\qquad\operatorname{Var}B=d/36,
 \qquad\operatorname{Cov}(A,B)=0.
\]

The limiting centered quadratic rows are \texttt{12\ A+(4/3)B} and
\texttt{30\ A+2B}; the nonconstant Casimir term is of order
\texttt{beta\^{}-1} on this exactly flat canonical chart. Hence their
covariance is

\[
 d\begin{pmatrix}1300/81&1082/27\\1082/27&901/9\end{pmatrix},
 \qquad \det=\frac{64}{81}d^2>0.
 \tag{31}
\]

This is a normalization and rank check only. It is not used in place of
the nonlinear good-bath proof (24)--(27), and is not assigned to general
actual baths.

\hypertarget{extensive-residuals-survive-every-complete-positive-observation}{%
\subsection{Extensive residuals survive every complete positive
observation}\label{extensive-residuals-survive-every-complete-positive-observation}}

Let \texttt{C\_a} stand for either actual dimensionless row
\texttt{h\_a(f)} or \texttt{R\_a(f)}, and set
\texttt{M\_a=E(C\_a\textbar{}B\_a)} and \texttt{epsilon\_a=C\_a-M\_a}.
Equations (4) and (7) prove

\[
 \|\epsilon_a\|_2^2\le C_f a^{-4}.
 \tag{32}
\]

Take any completed \textbf{positive} joint realization whose
original-link marginal is exactly the Wilson law, including auxiliary
variables, retained records and history labels. Let \texttt{Y} be any
actual later observation, and let \texttt{Q} be any orthogonal
projection on its output \texttt{L\^{}2} space, including a low-spectrum
projection if one has separately identified it. No screening relation
between \texttt{Y} and \texttt{B\_a} is required. Linearity and
conditional Jensen give

\[
 Q E(C_a\mid Y)=Q E(M_a\mid Y)+Q E(\epsilon_a\mid Y),
\]

\[
 \boxed{\|Q E(\epsilon_a\mid Y)\|_2^2\le C_f a^{-4}.}
 \tag{33}
\]

The same statement holds after centering, since centering is an
orthogonal contraction. Raw masses and independent scalar factors must
be restored if this probability identity is used inside an unnormalized
construction. No signed localized selected member is normalized into a
new probability law.

The triangle inequality shows that an extensive upper bound for
\texttt{Q\ E(C\_a\textbar{}Y)} is equivalent, up to the fixed additive
extensive residual, to an extensive upper bound for
\texttt{Q\ E(M\_a\textbar{}Y)}. This equivalence is in \texttt{L\^{}2}
norm. The difference of the \textbf{squared} norms can include the
uncontrolled cross term with \texttt{Q\ E(M\_a\textbar{}Y)} and is not
itself bounded by (33).

This provides a self-contained positive-observation bridge for the
actual ordinary row, based on the inputs from [ND] and the companion two-link
theorem, and on complete two-link stencils. It does not supply the
retained-mean estimate.

\newpage

\hypertarget{source-locators}{%
\section{Source locators}\label{source-locators}}

\begin{enumerate}
\def\labelenumi{\arabic{enumi}.}
\item
  \textbf{[W].} Wilson-source estimates, volumes 1--10
  (\texttt{lattice\_theory/source\_estimates} in this collection), cited by source page.
  Original source definitions, complete current/ordinary
  compensation, original scalar masses, and the open AE8/LM47, MT36 and
  S94 estimates (pp.~932--935, 781--784, 692 and 618--621).
\item
  \textbf{[ND].} \emph{Nonlinear dispersion and unequal spatial weights}
  (\texttt{conditional\_estimates/nonlinear\_dispersion\_and\_unequal\_spatial\_weights.tex}): nonlinear source dispersion and
  native covariance estimates. Source
  normalization: §1; original Wilson good-bath and complete one-link
  calculation: §§2.1--2.3; anti-concentration and non-tightness: §2;
  native scope: §§3--7.
\item
  \textbf{[UA].} \emph{Uniform averaging rows and finite-jet response}
  (\texttt{conditional\_estimates/uniform\_averaging\_rows\_and\_finite\_jet\_response.tex}): the incoming stored representative.
  Exact uniform flat averaging
  differential and counts; nonlinear and boundary remainders;
  same-body reserve and remaining incoming estimate.
\item
  \textbf{Companion notes.} The complete two-link theorem and flat jets are in
  \emph{A two-link obstruction to removing the physical H contact}
  (\texttt{local\_contacts/two\_link\_obstruction.tex}). The history-mean estimates
  are in [RS], \emph{Rare selector means and inherited coefficients}
  (\texttt{conditional\_estimates/rare\_selector\_means\_and\_inherited\_coefficients.tex}). The same-measure local exponential moments
  on odd Wilson tori are in \texttt{local\_contacts/odd\_torus\_transfer\_moment.tex}, with its inherited
  partition-power input and stated size/coupling regime.

\end{enumerate}
\end{document}
